\section{Transición a la segunda conferencia: de la memoria asociativa al Cognitive Map}

A partir de 00:42:28 cambia el ponente. El invitado previsto canceló a última hora, y Will Dorrell, con muy poco tiempo de preparación, dio la explicación escribiendo a mano en una tableta; por eso el formato visual pasa de un slide deck estándar a una página larga continua. El contenido, sin embargo, está estrechamente ligado a la primera mitad: Bricken pregunta «cómo una query similar recupera una memory», y Dorrell pregunta «cómo se recuerda la estructura de relaciones de la experiencia y cómo se reutiliza rápidamente con contenido nuevo». La primera pregunta enfatiza la geometría de direcciones; la segunda, la factorization entre estructura y contenido.

\subsection{Problema central: ¿puede reutilizarse la estructura de una experiencia secuencial?}

La primera página manuscrita reúne el hippocampus, el entorhinal system y los Transformers. La página siguiente usa un esquema de sequence/graph para ilustrar un problema general: los objetos y las recompensas que observa el agent pueden cambiar, pero la transition structure, la relative position o la relational rule pueden mantenerse sin cambios. Si el modelo memoriza desde cero el episode completo en cada entorno nuevo, no puede hacer few-shot generalize; si logra extraer la structure, puede transferir las reglas anteriores al contenido nuevo.

\lecturefigure{slide-41-cognitive-map-title.jpg}{Segunda conferencia: Hippocampus, entorhinal system y Transformers}{Estado didáctico manuscrito del video oficial V041; 00:43:54}

\lecturefigure{slide-42-one-problem-sequence-structure.jpg}{Un problema común: extraer una estructura reutilizable de la sequence experience}{Estado didáctico manuscrito del video oficial V042; 00:45:36}

\teachervoice{El moderador del curso explica que Will sustituye al invitado de manera provisional y presenta su formación en computational neuroscience en la Gatsby Unit. Acto seguido, Will describe su propio trabajo así: un modelo que originalmente estudiaba de forma independiente el hippocampal--entorhinal system y que después «resultó parecerse un poco a un Transformer». Esta frase enfatiza la convergent computation, no el diseño de un modelo del cerebro a imagen del Transformer.}

\begin{importantbox}{Definición mínima de structural generalization}
Dados entornos con distintas observation labels pero que comparten el mismo transition/relational graph, el sistema debe reutilizar, tras poca experiencia nueva, la estructura que ya tiene para hacer prediction, planning o inference. Lo que se transfiere es «cómo se conectan las relaciones», no la identidad sensorial de cada nodo.
\end{importantbox}

\subsection{Resumen de la sección}

La segunda conferencia lleva el problema de la memoria de la «recuperación por similitud» a la «reutilización de estructura». Un mapa cognitivo no solo debe recordar lo que ocurrió, sino también separar el what del where/how-related, para que los objetos nuevos puedan vincularse con rapidez a un grafo de relaciones conocido.
